\documentclass[../../main.tex]{subfiles}% 注意这里的文件路径不能用 ./main.tex ,否则用latexmk编译子文件会报错
\graphicspath{{\subfix{./image/}}{\subfix{../../image/}}} % 指定图片引用路径,后续可以直接使用图片文件名
% 如果图片存放在上级目录,则使用 ../../image/ 作为备用路径

% 例如:
% \begin{figure}[H]
% \centering
% \includegraphics[width=0.3\textwidth]{图.png}
% \caption{}
% \label{figure:图}
% \end{figure}
% 注意:上述\label{}一定要放在\caption{}之后,否则引用图片序号会只会显示??.

\begin{document}

\section{绝对连续函数与微积分基本定理}

\begin{lemma}\label{lemma:实变函数--引理5.9}
设 $f(x)$ 在 $[a,b]$ 上几乎处处可微, 且 $f'(x)=0$, a.e. $x\in[a,b]$. 若 $f(x)$ 在 $[a,b]$ 上不是常数函数, 则必存在 $\varepsilon>0$, 使得对任意的 $\delta>0$, $[a,b]$ 内存在有限个互不相交的区间:
\[
(x_1,y_1),\ (x_2,y_2),\ \cdots,\ (x_n,y_n),
\]
其长度的总和小于 $\delta$, 满足
\[
\sum_{i=1}^n |f(y_i)-f(x_i)|>\varepsilon.
\]
\end{lemma}

\begin{proof}
因为 $f(x)$ 不是常数, 所以不妨设存在 $c\in(a,b)$ 使得 $f(a)\neq f(c)$. 作点集 $E_c = \{x\in(a,b): f'(x)=0\}$. 对于任意的 $x\in E_c$, 由于 $f'(x)=0$, 故对任意的 $r>0$, 只要 $h$ 充分小, 且 $[x,x+h]\subset(a,c)$, 就有
\[
|f(x+h)-f(x)| < rh.
\]
于是 ($r$ 固定) 如此之区间 $[x,x+h]$ 的全体就构成 $E_c$ 的一个 Vitali 覆盖. 根据 Vitali 覆盖定理, 可知对任意 $\delta>0$, 存在互不相交的区间组:
\[
[x_1,x_1+h_1],\ [x_2,x_2+h_2],\ \cdots,\ [x_n,x_n+h_n],
\]
使得
\[
m\left(E_c\setminus\bigcup_{i=1}^n [x_i,x_i+h_i]\right) = m\left([a,c)\setminus\bigcup_{i=1}^n [x_i,x_i+h_i)\right) < \delta.
\]
不妨设这些区间的端点可排列为
\[
a = x_0 < x_1 < x_1+h_1 < x_2 < x_2+h_2 < \cdots < x_n+h_n < x_{n+1} = c.
\]
令 $h_0=0$, 并选 $\varepsilon>0$, 满足 $2\varepsilon < |f(c)-f(a)|$, 我们有
\begin{align*}
2\varepsilon &< |f(c)-f(a)| 
\leqslant \sum_{i=0}^n |f(x_{i+1})-f(x_i+h_i)| + \sum_{i=1}^n |f(x_i+h_i)-f(x_i)| \\
&\leqslant \sum_{i=0}^n |f(x_{i+1})-f(x_i+h_i)| + r\sum_{i=1}^n h_i \\
&\leqslant \sum_{i=0}^n |f(x_{i+1})-f(x_i+h_i)| + r(b-a).
\end{align*}
因为 $r$ 是任意的, 可先使得 $r(b-a)<\varepsilon$, 所以得到
\[
\varepsilon < \sum_{i=0}^n |f(x_{i+1})-f(x_i+h_i)|,
\]
而且
\[
\sum_{i=1}^n (x_{i+1}-(x_i+h_i)) = m\left(E_c - \sum_{i=1}^n (x_i,x_i+h_i)\right) < \delta.
\]
\end{proof}

\begin{definition}[绝对连续函数]\label{definecolor:绝对连续函数的定义}
设 $f(x)$ 是 $[a,b]$ 上的实值函数. 若对任给 $\varepsilon>0$, 存在 $\delta>0$, 使得当 $[a,b]$ 中任意有限个互不相交的开区间 $(x_i,y_i)$ ($i=1,2,\cdots,n$) 满足 $\sum_{i=1}^n (y_i-x_i) < \delta$ 时, 有
\[
\sum_{i=1}^n |f(y_i)-f(x_i)| < \varepsilon,
\]
则称 $f(x)$ 是 $[a,b]$ 上的\textbf{绝对连续函数}, 其全体记为 $AC([a,b])$.
\end{definition}
\begin{remark}
显然, 我们有下述简单事实:
\begin{enumerate}[(1)]
\item 绝对连续函数一定是连续函数;

\item 在 $[a,b]$ 上绝对连续函数的全体$AC([a,b])$构成一个线性空间.
\end{enumerate}
\end{remark}

\begin{proposition}\label{proposition:绝对连续函数对至多可数个开区间也成立}
设$f\in AC[a,b],a,b\in \mathbb{R}$,则对$\forall \varepsilon >0$,存在$\delta>0$,使得对任意可数个互不相交的开区间$\{(a_k,b_k)\}_{k=1}^{\infty}$满足$\sum\limits_{k=1}^{\infty}(b_k-a_k)<\delta$,有
\begin{align*}
\sum\limits_{k=1}^{\infty}\left|f(b_k)-f(a_k)\right|\leqslant \varepsilon .
\end{align*}
\end{proposition}
\begin{proof}
由$f\in AC[a,b]$知,对任给 $\varepsilon>0$, 存在 $\delta>0$, 使得当 $[a,b]$ 中任意有限个互不相交的开区间 $(x_i,y_i)$ ($i=1,2,\cdots,n$) 满足 $\sum_{i=1}^n (y_i-x_i) < \delta$ 时, 有
\begin{align}\label{eq::HJOIWIJOIJOIJOIJO23423242311}
\sum_{i=1}^n |f(y_i)-f(x_i)| < \varepsilon.
\end{align}
设可数个互不相交的开区间$\{(a_k,b_k)\}_{k=1}^{\infty}$满足$\sum\limits_{k=1}^{\infty}(b_k-a_k)<\delta$,则对$\forall N \in \mathbb{N}$,有
\begin{align*}
\sum\limits_{k=1}^{N}(b_k-a_k)<\delta.
\end{align*}
由\eqref{eq::HJOIWIJOIJOIJOIJO23423242311}式知
\begin{align*}
\sum_{k=1}^N |f(b_k)-f(a_k)| < \varepsilon.
\end{align*}
令$N\to \infty$得
\begin{align*}
\sum\limits_{k=1}^{\infty}\left|f(b_k)-f(a_k)\right|\leqslant \varepsilon .
\end{align*}
\end{proof}

\begin{proposition}\label{proposition:Lipschitz连续一定绝对连续}
若函数 $f(x)$ 在 $[a,b]$ 上满足 Lipschitz 条件:
\[
|f(x)-f(y)| \leqslant M|x-y|,\quad x,y\in[a,b],
\]
则 $f(x)$ 是 $[a,b]$ 上的绝对连续函数. 
\end{proposition}
\begin{proof}
事实上, 因为
\[
\sum_{i=1}^n |f(x_i)-f(y_i)| \leqslant M\sum_{i=1}^n |y_i-x_i| < \delta M,
\]
所以对于任意的 $\varepsilon>0$, 只需取 $\delta<\frac{\varepsilon}{M}$立即可知结论成立.
\end{proof}

\begin{theorem}
若 $f\in L([a,b])$, 则其不定积分
\[
F(x) = \int_a^x f(t)\mathrm{d}t
\]
是 $[a,b]$ 上的绝对连续函数.
\end{theorem}

\begin{proof}
对于任意的 $\varepsilon>0$, 因为 $f\in L([a,b])$, 所以存在 $\delta>0$, 当 $e\subset[a,b]$ 且 $m(e)<\delta$ 时, 有
\[
\int_e |f(x)|\mathrm{d}x < \varepsilon.
\]
现在, 对于 $[a,b]$ 中任意有限个互不相交的区间:
\[
(x_1,y_1),\ (x_2,y_2),\ \cdots,\ (x_n,y_n),
\]
当其长度总和小于 $\delta$ 时, 我们有
\begin{align*}
\sum_{i=1}^n |F(y_i)-F(x_i)| = \sum_{i=1}^n \left|\int_{x_i}^{y_i} f(x)\mathrm{d}x\right| \leqslant \sum_{i=1}^n \int_{x_i}^{y_i} |f(x)|\mathrm{d}x 
= \int_{\bigcup_{i=1}^n (x_i,y_i)} |f(x)|\mathrm{d}x < \varepsilon.
\end{align*}
\end{proof}

\begin{theorem}\label{theorem:AC一定BV}
若 $f(x)$ 是 $[a,b]$ 上的绝对连续函数, 则 $f(x)$ 是 $[a,b]$ 上的有界变差函数.
\end{theorem}
\begin{proof}
在函数绝对连续的定义中, 取 $\varepsilon=1$, 可知存在 $\delta>0$, 当 $[a,b]$ 中的任意有限个互不相交的开区间 $(x_i,y_i)$ ($i=1,2,\cdots,n$) 满足
\[
\sum_{i=1}^n (y_i-x_i) < \delta
\]
时, 必有
\[
\sum_{i=1}^n |f(y_i)-f(x_i)| < 1.
\]
作分划 $\Delta: a=c_0 < c_1 < \cdots < c_n = b$, 使得
\[
c_{k+1}-c_k < \delta,\quad k=0,1,\cdots,n-1,
\]
从而有
\[
\bigvee_{c_k}^{c_{k+1}}(f) \leqslant 1,\quad k=0,1,\cdots,n-1,
\]
故得
\[
\bigvee_a^b(f) \leqslant n.
\]
\end{proof}

\begin{corollary}\label{corollary:绝对连续函数几乎处处可微且导数可积}
若 $f(x)$ 是 $[a,b]$ 上的绝对连续函数, 则 $f(x)$ 在 $[a,b]$ 上是几乎处处可微的, 且 $f'(x)$ 是 $[a,b]$ 上的可积函数.
\end{corollary}
\begin{proof}
由\refthe{theorem:AC一定BV}和\refpro{proposition:有界变差函数必几乎处处可微}立得.
\end{proof}

\begin{theorem}
若 $f(x)$ 是 $[a,b]$ 上的绝对连续函数, 且 $f'(x)=0$, a.e. $x\in[a,b]$, 则 $f(x)$ 在 $[a,b]$ 上等于一个常数.
\end{theorem}
\begin{proof}
可由\reflem{lemma:实变函数--引理5.9}以及\refdef{definecolor:绝对连续函数的定义}直接推得.
\end{proof}


\begin{theorem}[微积分基本定理]\label{theorem:实变函数--微积分基本定理}
若 $f(x)$ 是 $[a,b]$ 上的绝对连续函数, 则
\[
f(x)-f(a) = \int_a^x f'(t)\mathrm{d}t,\quad x\in[a,b]. \tag{5.11}
\]
\end{theorem}

\begin{proof}
令
\[
g(x) = \int_a^x f'(t)\mathrm{d}t,\quad x\in[a,b],
\]
则 $g(x)$ 是 $[a,b]$ 上的绝对连续函数, 且有
\[
g'(x) = f'(x),\quad \text{a.e. } x\in[a,b],\quad g(a)=0.
\]
再令 $h(x)=f(x)-g(x)$, 则 $h(x)$ 也是绝对连续的, 且有
\[
h'(x) = f'(x)-g'(x) = 0,\quad \text{a.e. } x\in[a,b].
\]
从而可知 $h(x)$ 在 $[a,b]$ 上是一个常数, 而 $h(a)=f(a)$, 于是得
\begin{align*}
f(x) = h(x)+g(x) 
= f(a) + \int_a^x f'(t)\mathrm{d}t,\quad x\in[a,b].
\end{align*}
\end{proof}

\begin{corollary}[Lebesgue微积分基本定理]\label{corollary:Lebesgue微积分基本定理}
一个定义在 $[a,b]$ 上的函数 $f(x)$ 具有形式
\[
f(x) = f(a) + \int_a^x g(t)\mathrm{d}t,\quad g(t)\in L([a,b])
\]
的充分必要条件为 $f(x)$ 是 $[a,b]$ 上的绝对连续函数. 此时, 我们有
\begin{align*}
g(x)=f' (x),\quad \mathrm{a}.\mathrm{e}.\,x\in [a,b]
\end{align*}
\end{corollary}

\begin{proposition}
设 $g_k(x)$ ($k=1,2,\cdots$) 是 $[a,b]$ 上的绝对连续函数. 若
\begin{enumerate}[(1)]
\item 存在 $c$, $a\leqslant c\leqslant b$, 使得级数 $\sum_{k=1}^\infty g_k(c)$ 收敛;
\item $\sum_{k=1}^\infty \int_a^b |g_k'(x)|\mathrm{d}x < +\infty$,
\end{enumerate}
则级数 $\sum_{k=1}^\infty g_k(x)$ 在 $[a,b]$ 上是收敛的. 设其极限函数为 $g(x)$, 则 $g(x)$ 是 $[a,b]$ 上的绝对连续函数, 且有
\[
g'(x) = \sum_{k=1}^\infty g_k'(x),\quad \text{a.e. } x\in[a,b].
\]
\end{proposition}
\begin{proof}
由条件(ii)和\hyperref[theorem:非负可测函数的逐项积分定理]{非负可测函数的逐项积分定理}可知 函数
\[
G(x) = \sum_{k=1}^\infty g_k'(x),\quad x\in[a,b]
\]
是 $[a,b]$ 上的可积函数, 且有
\[
\lim_{n\to\infty} \int_c^x \sum_{k=1}^n g_k'(t)\mathrm{d}t = \int_c^x G(t)\mathrm{d}t.
\]
因为每个 $g_k(x)$ 都是绝对连续函数, 所以有
\[
g_k(x) = \int_c^x g_k'(t)\mathrm{d}t + g_k(c),\quad x\in[a,b],
\]
从而可知
\[
\sum_{k=1}^n g_k(x) = \int_c^x \sum_{k=1}^n g_k'(t)\mathrm{d}t + \sum_{k=1}^n g_k(c),\quad n=1,2,\cdots.
\]
现在令 $n\to\infty$, 即得
\[
\sum_{k=1}^\infty g_k(x) = \int_c^x G(t)\mathrm{d}t + \sum_{k=1}^\infty g_k(c),\quad x\in[a,b].
\]
若令上式左端为 $g(x)$, 则有
\[
g(x) = \int_c^x G(t)\mathrm{d}t + \sum_{k=1}^\infty g_k(c).
\]
由此可知 $g(x)$ 是 $[a,b]$ 上的绝对连续函数, 且有
\[
g'(x) = \sum_{k=1}^\infty g_k'(x),\quad \text{a.e. } x\in[a,b].
\]
\end{proof}

\begin{example}
证明不可能有 $[0,1]$ 中的可测集 $E$, 使得对一切 $x\in[0,1]$ 都有
\begin{align*}
m([0,x]\cap E)=\frac{1}{2}x.
\end{align*}
\end{example}
\begin{proof}
反证,设存在这样的可测集$E\subseteq [0,1]$,使得
\begin{align*}
m\left( [0,x] \cap E \right) =\frac{x}{2},\quad \forall x\in [0,1].
\end{align*}
记$f\left( x \right) =m\left( [0,x] \cap E \right) =\frac{x}{2}$,则
\begin{align*}
f\left( x \right) =m\left( [0,x] \cap E \right) =\int_0^x{\chi _E\left( t \right) \mathrm{d}t}.
\end{align*}
由\hyperref[corollary:Lebesgue微积分基本定理]{Lebesgue微积分基本定理}知
\begin{align*}
\chi _E\left( t \right) =f^\prime \left( x \right) =\frac{\mathrm{d}}{\mathrm{d}x}\left( \frac{x}{2} \right) =\frac{1}{2},\quad \mathrm{a}.\mathrm{e}.\,\,x\in [0,1].
\end{align*}
显然矛盾!
\end{proof}

\begin{example}[无处单调的绝对连续函数]
在 $[0,1]$ 中作点集 $E$, 使得 $[0,1]$ 中任一区间 $I$ 都有 (见\refpro{proposition:正测集与区间交测度小于区间})
\[
m(I\cap E) > 0,\quad m(I\cap E^c) > 0,
\]
并作 $[0,1]$ 上的绝对连续函数
\[
f(x) = \int_0^x [\chi_E(t) - \chi_{E^c}(t)]\mathrm{d}t.
\]
因此, 对 $[0,1]$ 中任一区间 $I$, 存在 $x_1\in I\cap E$, $x_2\in I\cap E^c$, 使得
\begin{align*}
f'(x_1) &= \chi_E(x_1) - \chi_{E^c}(x_1) = 1 > 0, \\
f'(x_2) &= \chi_E(x_2) - \chi_{E^c}(x_2) = -1 < 0.
\end{align*}
这说明 $f(x)$ 在区间 $I$ 上不是单调函数, 即得所证.
\end{example}

\begin{example}
若 $f(x)$ 在 $[a,b]$ 上绝对连续, 则弧长
\[
l_f = \int_a^b \sqrt{1+(f'(x))^2}\mathrm{d}x.
\]
\end{example}
\begin{proof}
实际上, 由\refpro{proposition:实变函数--弧长公式1}可知, 只需指出
\[
l_f \leqslant \int_a^b \sqrt{1+(f'(x))^2}\mathrm{d}x.
\]
为此, 对任给的 $\varepsilon>0$, 作 $[a,b]$ 的分划 $\Delta$, 使得 $l_\Delta(f)\geqslant l_f-\varepsilon$, 又作 $g\in C([a,b])$, 使得
\[
\int_a^b |f'(x)-g(x)|\mathrm{d}x < \varepsilon.
\]
现在令 $G(x) = \int_a^x g(t)\mathrm{d}t$, 易知 $|l_\Delta(f)-l_\Delta(G)| < \varepsilon$. 注意到
\begin{align*}
|\sqrt{1+\alpha^2}-\sqrt{1+\beta^2}| &\leqslant |\alpha-\beta|, \\
\left|\int_a^b \sqrt{1+(f'(x))^2}\mathrm{d}x - \int_a^b \sqrt{1+(g(x))^2}\mathrm{d}x\right|
&\leqslant \int_a^b |f'(x)-g(x)|\mathrm{d}x < \varepsilon.
\end{align*}
由此可知
\begin{align*}
l_\Delta(f) < l_\Delta(G) + \varepsilon \leqslant \int_a^b \sqrt{1+(G'(x))^2}\mathrm{d}x + \varepsilon 
= \int_a^b \sqrt{1+(g(x))^2}\mathrm{d}x + \varepsilon \leqslant \int_a^b \sqrt{1+(f'(x))^2}\mathrm{d}x + 2\varepsilon.
\end{align*}
注意到 $l_\Delta(f) > l_f-\varepsilon$, 最后导出
\[
l_f \leqslant \int_a^b \sqrt{1+(f'(x))^2}\mathrm{d}x + 3\varepsilon.
\]
由此即得所证.
\end{proof}

\begin{example}
设 $f\in L([c,d])$, $c<a<b<d$. 若有
\[
\int_a^b |f(x+h)-f(x)|\mathrm{d}x = O(|h|),\quad h\to 0,
\]
则 $f(x)=g(x)$, a.e. $x\in[a,b]$, 其中 $g\in  BV([a,b])$.
\end{example}

\begin{proof}
作 $[a,b]$ 上的函数列
\[
f_n(x) = n\int_x^{x+\frac{1}{n}} f(t)\mathrm{d}t\quad (n=1,2,\cdots),
\]
我们有
\begin{align*}
\int_a^b |f_n(x+h)-f_n(x)|\mathrm{d}x
&= n\int_a^b \left|\int_0^{\frac{1}{n}} (f(x+h+t)-f(x+t))\mathrm{d}t\right|\mathrm{d}x \\
&\leqslant n\int_0^{\frac{1}{n}} \left(\int_a^b |f(x+h+t)-f(x+t)|\mathrm{d}x\right)\mathrm{d}t \\
&= O(|h|),\quad h\to 0 \quad (\text{对 } n \text{ 一致}).
\end{align*}
由于 $f_n'\in L([a,b])$, 而且存在 $M>0$, 使得
\[
\int_a^b |f_n'(x)|\mathrm{d}x \leqslant \varliminf_{h\to 0} \int_a^b \left|\frac{f_n(x+h)-f_n(x)}{h}\right|\mathrm{d}x \leqslant M,
\]
故对 $[a,b]$ 的任一分割
\[
\Delta: a = x_0 < x_1 < \cdots < x_m = b,
\]
以及一切 $n$, 可得
\begin{align*}
\sum_{i=1}^m |f_n(x_i)-f_n(x_{i-1})| = \sum_{i=1}^m \left|\int_{x_{i-1}}^{x_i} f_n'(x)\mathrm{d}x\right| 
\leqslant \sum_{i=1}^m \int_{x_{i-1}}^{x_i} |f_n'(x)|\mathrm{d}x = \int_a^b |f_n'(x)|\mathrm{d}x \leqslant M.
\end{align*}
由此可知
\[
\bigvee_a^b(f_n) \leqslant M\quad (n=1,2,\cdots).
\]
注意到存在子列, 不妨仍记为 $\{f_n(x)\}$, 它在 $[a,b]\setminus Z$ ($m(Z)=0$) 上点收敛于 $f(x)$, 因此对分点不属于 $Z$ 的 $[a,b]$ 分割, 由
\[
\sum_{i=1}^m |f_n(x_i)-f_n(x_{i-1})| \leqslant M
\]
可得 (令 $n\to\infty$)
\[
\sum_{i=1}^m |f(x_i)-f(x_{i-1})| \leqslant M.
\]
由此即得所证.
\end{proof}

\begin{example}
设有函数 $f(y)=y^{1/3}$, $y\in[-1,1]$, 以及
\[
g(x) =
\begin{cases}
x^3\cos^3(\frac{\pi}{x}), & 0<x\leqslant 1, \\
0, & x=0,
\end{cases}
\]
易知 $f(y)$ 是 $[-1,1]$ 上的绝对连续函数, $g(x)$ 是 $[0,1]$ 上的绝对连续函数. 然而, 我们有
\[
F(x) = f(g(x)) =
\begin{cases}
x\cos\frac{\pi}{x}, & 0<x\leqslant 1, \\
0, & x=0,
\end{cases},
\quad \bigvee_0^1(F) = +\infty.
\]
\end{example}
\begin{remark}
这个例题表明,两个绝对连续函数的复合函数不一定是绝对连续的.
\end{remark} 

\begin{example}
在 $[0,1]$ 上作绝对连续函数列
\[
f_n(x) =
\begin{cases}
0, & 0\leqslant x\leqslant \frac{1}{n}, \\
x\sin\frac{\pi}{x}, & \frac{1}{n}<x\leqslant 1,
\end{cases}
\]
易知 $f_n(x)$ 在 $[0,1]$ 上一致收敛于 $f(x)$:
\[
f(x) =
\begin{cases}
0, & x=0, \\
x\sin\frac{\pi}{x}, & 0<x\leqslant 1,
\end{cases}
\]
而 $f(x)$ 在 $[0,1]$ 上不是有界变差的.
\end{example}
\begin{remark}
这个例题表明,绝对连续函数列在一致收敛运算下不封闭.
\end{remark}

\begin{example}
设
\[
f(x) =
\begin{cases}
x^2\sin\frac{1}{x}, & 0<x\leqslant 1, \\
0, & x=0.
\end{cases},
\]
则$f(x)$ 在 $[0,1]$ 上是绝对连续函数, 但 $|f(x)|^p$ ($0<p<1$) 在 $[0,1]$ 上不是绝对连续函数.
\end{example}
\begin{remark}
这个例题表明,$f(x)$ 在 $[0,1]$ 上是绝对连续函数, 但 $|f(x)|^p$ ($0<p<1$) 在 $[0,1]$ 上可以不是绝对连续函数.
\end{remark}

\begin{example}
设 $f(x)$ 在 $\mathbb{R}$ 上可微, 且有 $|f'(x)|\leqslant M$ ($x\in\mathbb{R}$). 若点集 $\{x\in\mathbb{R}: f'(x)>0\}$, $\{x\in\mathbb{R}: f'(x)<0\}$ 在 $\mathbb{R}$ 中稠密, 则对任意的 $[a,b]\subset\mathbb{R}$, $f'\notin R([a,b])$.
\end{example}
\begin{proof}
用反证法. 假定 $f'\in R([a,b])$, 则 $f'(x)$ 在 $[a,b]$ 上几乎处处连续, 且 $f\in  AC([a,b])$. 现在设 $x=x_0$ 是 $f'(x)$ 的连续点, 则 $f'(x_0)=0$. 这是因为如果 $f'(x_0)>0$, 那么由题设知存在 $\{x_n\}: x_n\to x_0$ ($n\to x_0$), $f'(x_n)<0$, 且有 $0\geqslant \lim_{n\to\infty} f'(x_n)=f'(x_0)$, 导致矛盾. 从而可得 $f'(x)=0$, a.e. $x\in[a,b]$, 这又导致 $f(x)=$ 常数 ($x\in[a,b]$), 矛盾. 证毕.
\end{proof}

\begin{proposition}\label{proposition:绝对连续函数零测集的像仍零测}
设 $f\in  AC([a,b])$, 则
\begin{enumerate}[(1)]
\item 对 $[a,b]$ 中的零测集 $Z$, 必有 $m(f(Z))=0$.

\item 若$E$ 是 $[a,b]$ 中的可测集, 则 $f(E)$ 是可测集.
\end{enumerate}
\end{proposition}
\begin{proof}
\begin{enumerate}[(1)]
\item 设$Z\subseteq [a,b]$是零测集，则由外测度定义知，对任意固定的$n\in \mathbb{N}$，存在开区间列$\{I_{ni}\}_{i=1}^{\infty}\subseteq [a,b]$，使得
\begin{align*}
Z\subseteq \bigcup\limits_{i=1}^{\infty}I_{ni},\quad \sum\limits_{i=1}^{\infty}|I_{ni}|<\frac{1}{n}.
\end{align*}
令$J_{ni}=I_{ni}\setminus \bigcup\limits_{j=1}^{i-1}I_{nj}\ (\forall i\in \mathbb{N})$，则
\begin{align}
Z\subseteq \bigcup\limits_{i=1}^{\infty}I_{ni}=\bigcup\limits_{i=1}^{\infty}J_{ni}\subseteq \bigcup\limits_{i=1}^{\infty}\overline{J_{ni}},\quad \sum\limits_{i=1}^{\infty}|J_{ni}|\leqslant \sum\limits_{i=1}^{\infty}|I_{ni}|<\frac{1}{n},\quad J_{ni}\cap J_{nj}=\varnothing\ (\forall i\neq j). \label{eq::::--234823urjreosew1}
\end{align}
因为$f\in AC[a,b]$，所以$f\in C[a,b]$，进而对$\forall i\in \mathbb{N}$，$f$在有界闭集$\overline{J_{ni}}$上存在最小值$f(a_{ni})$和最大值$f(b_{ni})$，其中$a_{ni},b_{ni}\in \overline{J_{ni}}$。于是再结合连续函数介值定理知
\begin{align}
f\left( \overline{J_{ni}} \right) =\left[ f\left( a_{ni} \right) ,f\left( b_{ni} \right) \right],\quad \forall i\in \mathbb{N}. \label{eq321}
\end{align}
记
\begin{align*}
K_{ni}\triangleq \left( \min \left( a_{ni},b_{ni} \right) ,\max \left( a_{ni},b_{ni} \right) \right),\quad \forall i\in \mathbb{N}.
\end{align*}
由连续函数介值定理和\eqref{eq321}式知
\begin{align*}
f\left( K_{ni} \right) =\left[ f\left( a_{ni} \right) ,f\left( b_{ni} \right) \right] =f\left( \overline{J_{ni}} \right),\quad \forall i\in \mathbb{N}.
\end{align*}
再结合\eqref{eq::::--234823urjreosew1}式知，可数个开区间$\{(a_i,b_i)\}_{i=1}^{\infty}\subseteq [a,b]$满足
\begin{gather}
\sum\limits_{i=1}^{\infty}|K_{ni}|\leqslant \sum\limits_{i=1}^{\infty}|J_{ni}|<\frac{1}{n},\quad K_{ni}\cap K_{nj}=\varnothing\ (\forall i\neq j),\label{eq::::--234823urjreosew2-1}\\
m^*\left( f\left( \overline{J_{ni}} \right) \right) =m^*\left( f\left( K_{ni} \right) \right) =f\left( b_{ni} \right) -f\left( a_{ni} \right)\ (\forall i\in \mathbb{N}). \label{eq::::--234823urjreosew2}
\end{gather}
由$f\in AC[a,b]$和\refpro{proposition:绝对连续函数对至多可数个开区间也成立}知，对$\forall \varepsilon >0$，$\exists \delta >0$，使得当任意可数个互不相交开区间$\{(a_i,b_i)\}_{i=1}^{\infty}\subseteq [a,b]$满足$\sum\limits_{i=1}^{\infty}(b_i-a_i)<\delta$时，有
\begin{align}
\sum\limits_{i=1}^{\infty}|f(b_i)-f(a_i)|<\varepsilon. \label{eq::::--234823urjreosew3}
\end{align}
当$n>\frac{1}{\delta}$时，由\eqref{eq::::--234823urjreosew2-1}式知$\sum\limits_{i=1}^{\infty}|K_{ni}|<\frac{1}{n}$，故由\eqref{eq::::--234823urjreosew2}\eqref{eq::::--234823urjreosew3}式知
\begin{align*}
\sum\limits_{i=1}^{\infty}m^*\left( f\left( K_{ni} \right) \right) =\sum\limits_{i=1}^{\infty}\left[ f\left( b_{ni} \right) -f\left( a_{ni} \right) \right] \leqslant \varepsilon.
\end{align*}
于是再结合\eqref{eq::::--234823urjreosew1}\eqref{eq::::--234823urjreosew2-1}\eqref{eq::::--234823urjreosew2}式以及\rrefpro{Set Theory-proposition:映射的基本性质}{Set Theory-proposition:映射的基本性质-2}知
\begin{align*}
m^*\left( f\left( Z \right) \right) \leqslant m^*\left( f\left( \bigcup\limits_{i=1}^{\infty}\overline{J_{ni}} \right) \right) =m^*\left( \bigcup\limits_{i=1}^{\infty}f\left( \overline{J_{ni}} \right) \right) \leqslant \sum\limits_{i=1}^{\infty}m^*\left( f\left( \overline{J_{ni}} \right) \right) =\sum\limits_{i=1}^{\infty}m^*\left( f\left( K_{ni} \right) \right) \leqslant \varepsilon.
\end{align*}
令$\varepsilon \rightarrow 0^+$得$m^*(f(Z))=0$，即$f(Z)$是零测集。

\item 若$E\subseteq [a,b]$是可测集，则由\rrefthe{theorem:Lebesgue可测集的充要条件}{theorem:Lebesgue可测集的充要条件-5}知，存在$F_{\sigma}$集$F\supseteq E$，使得
\begin{align*}
m(F\setminus E)=0.
\end{align*}
由(1)知$f(F\setminus E)$是零测集。从而由\rrefpro{Set Theory-proposition:映射的基本性质}{Set Theory-proposition:映射的基本性质-5}
知
\begin{align*}
f(F)\setminus f(E)\subseteq f(F\setminus E) \Longrightarrow m^*\left( f(F)\setminus f(E) \right) \leqslant m^*\left( f(F\setminus E) \right) =0.
\end{align*}
故$f(F)\setminus f(E)$也是零测集。

由$F$是$F_{\sigma}$集知，可不妨设$F=\bigcup\limits_{n=1}^{\infty}F_n$，其中$F_n\ (\forall n\in \mathbb{N})$都是闭集。对$\forall n\in \mathbb{N}$，由$f\in AC[a,b]$知$f\in C[a,b]$，故对$\forall n\in \mathbb{N}$，$f$在$F_n$上存在最小值$f(c_n)$和最大值$f(d_n)$，其中$c_n,d_n\in F_n$。从而由连续函数介值定理知
\begin{align*}
f(F_n)=\left[ f(c_n),f(d_n) \right],\quad \forall n\in \mathbb{N}.
\end{align*}
因为$\mathbb{R}$上的闭区间都可测，所以$f(F_n)\ (\forall n\in \mathbb{N})$都可测，因此由\rrefpro{Set Theory-proposition:映射的基本性质}{Set Theory-proposition:映射的基本性质-2}知$f(F)=\bigcup\limits_{n=1}^{\infty}f(F_n)$也是可测集。于是再由$f(F)\setminus f(E)$是零测集知$f(E)=f(F)\setminus (f(F)\setminus f(E))$是可测集。
\end{enumerate}
\end{proof}

\begin{example}
设 $f\in C([a,b])\cap  BV([a,b])$. 若对 $[a,b]$ 中的任一零测集 $Z$, 必有 $m(f(Z))=0$, 则 $f\in  AC([a,b])$.
\end{example}
\begin{proof}
用反证法. 假定 $f\notin  AC([a,b])$, 则存在 $\varepsilon_0>0$, 使得对每个 $n\in\mathbb{N}$, 都有 $[a,b]$ 的子区间组
\[
I_{n,k} = [a_{n,k}, b_{n,k}] \quad (k=1,2,\cdots,j_n), \quad \sum_{k=1}^{j_n} (b_{n,k} - a_{n,k}) \leqslant \frac{1}{2^n},
\]
\[
\sum_{k=1}^{j_n} |f(b_{n,k}) - f(a_{n,k})| \geqslant \varepsilon_0.
\]
作点集 $E_i = \bigcup_{n=i}^{\infty} \bigcup_{k=1}^{j_n} I_{n,k}$ ($i\in\mathbb{N}$), 我们有
\[
m(E_i) \leqslant \sum_{n=i}^{\infty} \sum_{k=1}^{j_n} (b_{n,k} - a_{n,k}) \leqslant \sum_{n=i}^{\infty} \frac{1}{2^n} = \frac{1}{2^{i-1}}.
\]
由此可知 $m(E_i)\to 0$ ($i\to\infty$). 因为 $f\in C([a,b])$, 所以 $f(E_i)$ 是紧区间的并集. 又由 $[a,b]\supset E_i\supset E_{i+1}\supset\cdots$ 可得
\[
m(E) = 0 \quad \left(E = \bigcap_{i=1}^{\infty} E_i\right), \quad m(f(E)) = 0,
\]
从而根据 $f([a,b])\supset f(E_i)\supset f(E_{i+1})\supset\cdots$, 又有
\begin{align}
\lim_{i\to\infty} m(f(E_i)) = m(f(E)) = 0. \label{eq::923jri2jrijisofgzdx}
\end{align}
另一方面, 若 $f(x)$ 在 $[a,b]$ 上递增, 易知
\[
f(I_{i,k}) = [f(a_{i,k}), f(b_{i,k})], \quad m(f(E_i)) \geqslant \sum_{k=1}^{j_i} |f(b_{i,k}) - f(a_{i,k})| \geqslant \varepsilon_0.
\]
这与\eqref{eq::923jri2jrijisofgzdx}式矛盾！
\end{proof}











\end{document}